%% LyX 2.4.4 created this file.  For more info, see https://www.lyx.org/.
%% Do not edit unless you really know what you are doing.
\documentclass[english]{article}
\usepackage[T1]{fontenc}
\usepackage[latin9]{inputenc}
\usepackage{enumitem}
\usepackage{amsmath}
\usepackage{amsthm}
\usepackage{amssymb}
\usepackage{cancel}
\usepackage{geometry}
\geometry{verbose,tmargin=1in,bmargin=1in,lmargin=1in,rmargin=1in}

\makeatletter
%%%%%%%%%%%%%%%%%%%%%%%%%%%%%% Textclass specific LaTeX commands.
\numberwithin{equation}{section}
\numberwithin{figure}{section}
\newlength{\lyxlabelwidth}      % auxiliary length 

%%%%%%%%%%%%%%%%%%%%%%%%%%%%%% User specified LaTeX commands.
\usepackage{fancyhdr}
\usepackage{lastpage}
\pagestyle{fancy}
\usepackage{enumitem}
\usepackage{ifthen}
\everymath=\expandafter{\the\everymath\displaystyle}
%\DeclareMathSizes{10}{10}{10}{10}
\usepackage{multicol}

\usepackage{tikz}
\usetikzlibrary{patterns}
\usetikzlibrary{plotmarks}
\usepackage{pgfplots}
\pgfplotsset{compat=newest}

\usepackage{advdate}    % Advancing/saving dates
\usepackage{datetime}   % Dates formatting
\usepackage{datenumber} % Counters for dates
\usdate
% Added by lyx2lyx
\setlength{\parskip}{\smallskipamount}
\setlength{\parindent}{0pt}

\makeatother

\usepackage{babel}
\begin{document}
\renewcommand{\author}{Dr.\,James Ashe}

\newcommand{\classNum}{335}

\newcommand{\classSec}{A}

\newcommand{\examType}{Practice Midterm}

\renewcommand{\date}{\formatdate{22}{3}{2013}}

%%First page's header%%%%%%%%%%%%%%%%%%%%%%
\fancypagestyle{firststyle} %%header settings for first pagr
{
	\fancyhead[L]{MTH \classNum{}(\classSec{}) \author{}\\
	\textbf{\examType{}} \date{}}
	\fancyhead[C]{\textbf{Name:}}
	\setlength{\headsep}{14pt}
}
%%%%%%%%%%%%%%%%%%%%%%%%%%%%%%%%%%%%%%%%%%

%%Next page's header%%%%%%%%%%%%%%%%%%%%%%%
\setlist{nolistsep}
\lhead{\classNum{}(\classSec{})} 
\chead{\textbf{}} 
\cfoot{\thepage{}~of~\pageref{LastPage}} 
\setlength{\headsep}{5pt} 
%%%%%%%%%%%%%%%%%%%%%%%%%%%%%%%%%%%%%%%%%%%

%%Various settings; see the preamble for other settings%%%%%
%\setenumerate[0]{leftmargin=12pt,itemindent=0pt,labelwidth=0pt}
\setlist{topsep=.5pt}
\renewcommand{\labelenumi}{\textbf{\arabic{enumi}.}}
\renewcommand{\labelenumii}{\textbf{\alph{enumii}.}}
\thispagestyle{firststyle}
%%%%%%%%%%%%%%%%%%%%%%%%%%%%%%%%%%%%%%%%%%

\newcommand{\Dspace}{.5cm}
\begin{enumerate}[resume]
\item Determine the following

\begin{enumerate}
\item $17\bmod2$\vspace{\Dspace}
\item $\mbox{lcm}(7^{4000}\cdot5^{2}\cdot2^{4},7^{4}\cdot5^{1000}\cdot2)$\vspace{\Dspace}
\item $17=x\cdot2-y\cdot15$. What is $\gcd(x,y)$?\vspace{\Dspace}
\end{enumerate}
\item Complete the following

\begin{enumerate}
\item Write the given set in tabular form: $A=\{x\mid x\in\mathbb{N}\mbox{ and }x\leq5\}$\vspace{\Dspace}\vspace{\Dspace}
\item Write the given set in set-builder form: $A=\left\{ \dots-2,0,2,4\right\} $\vspace{\Dspace}\vspace{\Dspace}
\end{enumerate}
\item Let $A=\left\{ 1,2,3\right\} $, $B=\left\{ 2,4,6\right\} $, $C=\left\{ 4,5,6\right\} $.

\begin{enumerate}
\item Find $A\cup C$\vspace{\Dspace}\vspace{\Dspace}
\item Find $A\cap B$\vspace{\Dspace}\vspace{\Dspace}
\item Find $A\cap C$\vspace{\Dspace}\vspace{\Dspace}
\item Find $\left(A\cup C\right)\cap B$\vspace{\Dspace}\vspace{\Dspace}
\end{enumerate}
\item Let $A$ and $B$ be two sets. Prove that $A\cap B\subseteq A$. \vfill
\item Let $A\neq\emptyset$ be a subset of $B$. Prove that $A\cup B=A$.\vfill\newpage
\item Give the definition of the following:

\begin{enumerate}
\item A \emph{function}\vspace{\Dspace}\vspace{\Dspace}\vspace{\Dspace}
\item An \emph{onto function}.\vspace{\Dspace}\vspace{\Dspace}\vspace{\Dspace}
\end{enumerate}
\item Consider the map $f:\mathbb{R}\times\mathbb{R}\rightarrow\mathbb{R}$
defined as $f(x,y)=xy$. Then $f$ is:

\begin{enumerate}
\item [(a)] an onto function.
\item [(b)] a 1-1 function.
\item [(c)] a bijection.
\item [(d)] none of the above.\vspace{\Dspace}
\end{enumerate}
\item Let $f:A\rightarrow B$ be a function, $A=\left\{ 1,2,3\right\} $,
and $B=\left\{ 1,2\right\} $. 

\begin{enumerate}
\item Can $f$ be 1-1? Explain.\vspace{\Dspace}\vspace{\Dspace}\vspace{\Dspace}
\item Show that $f$ can be onto. \vspace{\Dspace}\vspace{\Dspace}\vspace{\Dspace}\vspace{\Dspace}
\end{enumerate}
\item Define the following relation on the real number set: $x\mathcal{R}y\iff xy>0$.
Show the following:

\begin{enumerate}
\item $\mathcal{R}$ is \emph{not }reflexive.\vspace{\Dspace}\vspace{\Dspace}
\item $\mathcal{R}$ is transitive.\vspace{\Dspace}\vspace{\Dspace}\vspace{\Dspace}
\item $\mathcal{R}$ is \emph{not }an equivalence relation.\vspace{\Dspace}\vspace{\Dspace}
\end{enumerate}
\item Give an example of a relation on a set which is an equivalence relation.
\end{enumerate}

\end{document}
